% Notes of the Montreal seminar, 1970 (cote 129, typed part): archive cote 129, site page 54 (PDF page 55).
% typescript. Draft compared: none.
% AI-transcribed by Claude Opus 5.5 (claude-opus-5-5, adaptive thinking, effort high), Message Batches API, batch mtl129_full (msgbatch_01AGnDKHwy6ATCM7ehUa6263), 2026-11-13; not yet checked by eye.
\scanpage{129/54}
\mspage{\hand{II.26}}

pour tout \ $x \in M$ , car pour \ $u$ , $u' \in \mathrm{Hom}_{\mathcal{D}}(M, \mathbb{W}_{N+1}(S))$ , on a
$\phi_S(u + u')(T_x) = S_N(u(x), u'(x))$ \ et, toujours d'après la définition
de \ $\phi_S$ , on a, en posant \ $\overline{u} = \phi_S(u)$ , $\overline{u}' = \phi_S(u')$
\[ S_N(u(x),u'(x)) = S_N[\overline{u}(T_{V^N(x)}),\dots,\overline{u}(T_x) \;;\; \overline{u}'(T_{V^N(x)}),\dots,\overline{u}'(T_x)] \ . \]

De sorte que \ $\overline{u}$ \ est un morphisme de groupe si et seule-%
ment si pour tout \ $x \in M$ , on a
\[ \Delta(\overline{u}(T_x)) = S_N[\overline{u}(T_{V^N(x)}) \otimes 1,\dots,\overline{u}(T_x) \otimes 1;\ 1 \otimes \overline{u}(T_{V^N(x)}),\dots,1 \otimes \overline{u}(T_x)] \ . \]

Ce dernier élément n'est autre que la N+1$^{\text{ième}}$ coordonnée
\margin{$\times$ \ $\otimes$} dans \ $\mathbb{W}_{N+1}(\Gamma(\underline{\mathcal{O}}_G)$ \hand{(}$x$\hand{)} $\Gamma(\underline{\mathcal{O}}_G))$ \ de l'élément \ $u(x) \otimes 1 + 1 \otimes u(x)$ .

Le même raisonnement appliqué à \ $x$ , $V(x),\dots,V^N(x)$ \ montre
alors que \ $\overline{u}$ \ définit un morphisme de groupes \ $G \to \mathbb{E}(M)$ \ si et seulement
si \ $u(x)$ \ appartient au noyau de \ $\mathbb{W}_{N+1}(\Gamma(\underline{\mathcal{O}}_G)) \xrightarrow{\ \Delta^*\text{-}\mathrm{pr}_1^*\text{-}\mathrm{pr}_2^*\ } \mathbb{W}_{N+1}(\Gamma(\underline{\mathcal{O}}_G) \otimes \Gamma(\underline{\mathcal{O}}_G))$ ,
qui est lui-même contenu dans \ Prim(G) .

\margin{54}
